\section{The Names of the Angels}\label{app:names}

Scripture gives the proper names of three holy angels, Michael, Gabriel,
and Raphael, and the Church names no others. The Roman synod of 745 confessed
that it knew the names of no more than three: \latin{non plus quam trium
angelorum nomina cognoscimus, id est Michael, Gabriel, Raphael} (MGH
\emph{Conc.}\ II/1, ed.\ Werminghoff, 1906, p.~43); Charlemagne's \emph{Admonitio
generalis} of 789 forbade that unknown names be invented or used,
\latin{nisi illos quos habemus in auctoritate, id sunt Michahel,
Gabrihel, Rafahel} (c.~16, MGH \emph{Capit.}\ I, ed.\ Boretius, 1883, p.~55); and the
Directory of 2001 teaches that ``the practice of assigning names to the
Holy Angels should be discouraged, except in the cases of Gabriel,
Raphael and Michael whose names are contained in Holy Scripture''
(\emph{Directory on Popular Piety} 217). The acts are expounded in
Section~\ref{sec:devotion}.

The word ``angel'' is itself a name of office. Gregory teaches,
\latin{angelorum vocabulum, nomen est officii, non naturae}, and Isidore
repeats it, \latin{Angelorum autem vocabulum officii nomen est, non
naturae. Semper enim spiritus sunt, sed cum mittuntur, vocantur angeli}
(Gregory, \emph{Hom.\ in Evang.} 34.8; Isidore, \emph{Etym.}\ VII.5.2).
The proper names are names of office too. In the heavenly city, Gregory
says, the angels do not take names in order to be known, for the vision
of God makes all known; \latin{sed cum ad nos aliquid ministraturi
veniunt, apud nos etiam nomina a ministeriis trahunt}: when they come to
us on some ministry, they take their names among us from their
ministries (\emph{Hom.} 34.8).

\subsection*{The Three Names of Scripture}

\begin{fourstable}{Name}{Scripture (Vulgate)}{Meaning as the Fathers give it}{In the Roman liturgy}
Michael & Dan 10:13, 21; 12:1; Jude 9; Apoc 12:7 & \latin{Quis ut
Deus} (Gregory, \emph{Hom.} 34.9); \latin{qui interpretatur quis sicut
Deus} (Jerome, \emph{In Dan.}\ 8:16); \latin{Qui sicut Deus} (Isidore,
\emph{Etym.}\ VII.5.12) & 29 September; 8 May (pre-1955); Confiteor;
incense prayer; Offertory of the Requiem; Litany\\
Gabriel & Dan 8:16; 9:21; Luke 1:19, 26 & \latin{fortitudo Dei}
(Gregory, \emph{Hom.} 34.9); \latin{fortitudo, vel robustus Dei}
(Jerome, \emph{In Dan.}\ 8:16); \latin{fortitudo Dei} (Isidore,
\emph{Etym.}\ VII.5.10) & 24 March (pre-1955, 1962); 29 September
(postconciliar); the Angelus; Litany\\
Raphael & Tob 3:25; 5:5; 12:15 & \latin{medicina Dei} (Gregory,
\emph{Hom.} 34.9); \latin{curatio; vel medicina Dei} (Jerome, \emph{In
Dan.}\ 8:16); \latin{curatio vel medicina Dei} (Isidore, \emph{Etym.}\
VII.5.13) & 24 October (pre-1955, 1962); 29 September
(postconciliar); Litany\\
\end{fourstable}

The Fathers read each name as the sign of the angel's work. For Gregory,
whenever a work of wondrous power is to be done Michael is sent, so that
by the deed and the name it may be understood \latin{quia nullus potest
facere quod facere praevalet Deus}, that no one can do what God can do;
and the old enemy who in his pride desired to be like God, saying
\latin{similis ero Altissimo}, will at the end of the world be slain by
Michael and learn \latin{quia ad Dei similitudinem per superbiam nullus
exsurgat}. Gabriel, the Strength of God, is sent to Mary because he came
to announce the one who came humble to war against the powers of the air,
the \latin{Dominus fortis et potens, Dominus potens in praelio} of the
psalm. Raphael is the Medicine of God because he touched the eyes of
Tobias as a physician and wiped away the darkness of his blindness
(Gregory, \emph{Hom.} 34.9; read in the Breviary on 29 September,
Section~\ref{sec:liturgy}). Jerome, commenting on Daniel's vision, gives
the same three offices in other words: Gabriel, \latin{qui praepositus est
praeliis}, is sent where wars and the succession of kingdoms are at
issue; wherever healing is needed, \latin{Raphael mittitur}; and where
propitiation or expiation is needed, \latin{Michael dirigitur, qui
interpretatur quis sicut Deus} (Jerome, \emph{In Dan.}\ 8:16, PL~25).
Isidore gathers the whole into his \emph{Etymologies}: where divine power
or strength is shown, Gabriel is sent; when a work of wondrous power is
done, Michael; wherever there is need of healing, Raphael (\emph{Etym.}\
VII.5.10--14).

Scripture gives Raphael a second name in the book of Tobias, the name he
took in his disguise: ``I am Azarias the son of the great Ananias'' (Tob
5:18). The sermon the Breviary for England reads under Augustine's name
on Raphael's feast comments: ``The Angel hid his name from you; he said
not, I am the Angel Raphael, but, I am Azarias, the son of Ananias the
Great. In the first name that he gave thee he hid his high estate, lest
he should scare him that hired him'' (Bute, IV, p.~686).

\subsection*{The Titles of Michael in the Liturgy}

\begin{threestable}{Title}{Where the Church uses it}{Scripture behind it}
\latin{Princeps militiae Angelorum}; ``Prince of the heavenly host'' &
Responsory of 29 September (\emph{BR}); prayer after Low Mass (``O Prince
of the heavenly host'') & Dan 10:13, 21; 12:1; Apoc 12:7\\
\latin{Praepositus paradisi}, ``the Vice-Roy of Paradise'' & Antiphon and
responsory of 29 September (\emph{BR}; Bute, IV, p.~594) & ---\\
\latin{Salutis signifer}; \latin{signifer sanctus Michael} & Vespers hymn
of 29 September; Offertory of the Mass for the Dead & Apoc 12:7--9\\
\latin{Angelus pacis} & Lauds hymn of 29 September & ---\\
\latin{Dei nuntius pro animabus justis}; prince \latin{super omnes animas
suscipiendas} & Antiphons of 29 September & Luke 16:22\\
\latin{Princeps gloriosissime} & Magnificat antiphon of 29 September;
\emph{Raccolta} n.~289 & ---\\
``Prince of the Church of Jesus Christ'' & Versicle of the chaplet
(\emph{Raccolta} n.~291) & Dan 12:1\\
\end{threestable}

\subsection*{The Names of the Fallen in Scripture}

Scripture gives the fallen angels names that are titles of their revolt
and of their work. The Church uses them as Scripture does.

\begin{threestable}{Name}{Scripture (Douay)}{Reading}
Satan & ``Satan also was present among them'' (Job 1:6); ``that old
serpent, who is called the devil and Satan'' (Apoc 12:9) & The accuser:
``the accuser of our brethren is cast forth'' (Apoc 12:10)\\
The devil (\latin{diabolus}) & ``to be tempted by the devil'' (Matt
4:1); Apoc 12:9 & The tempter of Christ, ``who seduceth the whole world''
(Apoc 12:9)\\
The dragon, the old serpent & Apoc 12:7--9 & ``that old serpent,'' cast out by Michael and his angels (Apoc 12:9)\\
Beelzebub & ``Beelzebub the prince of the devils'' (Matt 12:24) & The
name the Pharisees used for the prince of the demons\\
Abaddon, Apollyon, \latin{Exterminans} & ``the angel of the bottomless
pit (whose name in Hebrew is Abaddon and in Greek Apollyon, in Latin
Exterminans)'' (Apoc 9:11) & The king over the locusts of the abyss
(Apoc 9:3--11)\\
Asmodeus & ``a devil named Asmodeus'' (Tob 3:8) & The demon bound by
Raphael (Tob 8:3)\\
Belial & ``what concord hath Christ with Belial?'' (2 Cor 6:15) & Set
by the Apostle over against Christ\\
The prince of this world & ``now shall the prince of this world be cast
out'' (John 12:31) & The dominion broken by the Cross
(Section~\ref{sec:assaults})\\
Lucifer & ``How art thou fallen from heaven, O Lucifer, who didst rise
in the morning?'' (Isa 14:12) & Applied by the Fathers to the devil's
fall: Origen, \emph{De principiis} I.5.5; Gregory, \emph{Hom.} 34.9 on
Isa 14:13--14\\
\end{threestable}

The prophet speaks to the king of Babylon, and the Fathers read in his
fall the fall of the first rebel. Origen argues from the name itself:
``Most evidently by these words is he shown to have fallen from heaven,
who formerly was Lucifer, and who used to arise in the morning. For if,
as some think, he was a nature of darkness, how is Lucifer said to have
existed before? Or how could he arise in the morning, who had in himself
nothing of the light?'' (Origen, \emph{De principiis} I.5.5, trans.\
Crombie, ANF~4). The name of the light-bearer is thus a witness to the
doctrine the Church defined against the dualists: the devil was created
good and became evil of himself (Lateran~IV, DS~800; Section~\ref{sec:faith};
Section~\ref{sec:fall}). Scripture adds that ``Satan himself transformeth
himself into an angel of light'' (2 Cor 11:14).

\subsection*{Names from the Apocryphal Books}

The fourth book of Esdras, which the Douay Bible of 1610 printed after
the canonical books ``because they are not receiued into the Canon of
Diuine Scriptures by the Catholique Church,'' names two angels: ``the
Angel answered me, that was sent to me, whose name was Vriel,'' and
``Ieremiel the Archangel answered to those things'' (4 Esdras 4:1, 36,
Douay 1610 appendix). Isidore received the first name into his
\emph{Etymologies}: \latin{Vriel interpretatur ignis Dei, sicut legimus
apparuisse ignem in rubo} (\emph{Etym.}\ VII.5.15). The first book of
Enoch gives seven ``names of the holy angels who watch'': Uriel,
Raphael, Raguel, Michael, Saraqâêl, Gabriel, and Remiel (1~En 20:1--8,
trans.\ Charles, 1917; Section~\ref{sec:before-christ}). The names that
stand only in these books are graded Apocryphal.

\subsection*{The Names the Roman Synod of 745 Rejected}

Aldebert's prayer, read before Pope Zachary's synod, invoked eight angels
by name: \latin{angelus Uriel, angelus Raguel, angelus Tubuel, angelus
Michael, angelus Adinus, angelus Tubuas, angelus Sabaoc, angelus Simiel}
(MGH \emph{Conc.}\ II/1, p.~42). The synod judged that of these only
Michael was an angel, and that under the names of angels Aldebert had
invoked demons, \latin{sub obtentu angelorum demonum nomina introduxit}
(ibid., p.~43).

\begin{threestable}{Name}{Other witness}{Disposition}
Uriel & 4 Esdras 4:1 (``Vriel''); Isidore, \emph{Etym.}\ VII.5.15 &
Rejected by the Roman synod of 745; Apocryphal\\
Raguel & 1~En 20:4 (Section~\ref{sec:before-christ}) & Rejected by the
Roman synod of 745; Apocryphal\\
Tubuel, Adinus, Tubuas, Sabaoc, Simiel & Known from Aldebert's prayer &
Rejected by the Roman synod of 745 as names of demons\\
Michael & Dan 10:13; Jude 9; Apoc 12:7 & The one angel's name in the
prayer; a name of Scripture\\
\end{threestable}

The Congregation for the Doctrine of the Faith applied the same rule in
1983 and 1992 to names of angels drawn from presumed private
revelations, which may be neither taught nor used in prayer (decree of
6 June 1992, AAS 84 [1992] 805--806; Section~\ref{sec:devotion}).
